
\documentclass[12pt]{article}
\usepackage{amsmath, amssymb, physics, geometry}
\usepackage{bm}
\geometry{margin=1in}
\title{\textbf{SAT.4D Covariant Field Equations}}
\author{SAT Framework Team}
\date{June 2025}

\begin{document}

\maketitle

\section*{1. Introduction}

This section presents the SAT.4D covariant field equations as they arise from the topological and statistical structure of 1D filaments embedded in a 4D manifold. All equations are expressed using emergent quantities derived from filament bundle configurations, with no a priori field insertions.

\section*{2. Scalar Field: Emergent Klein--Gordon Equation}

\subsection*{2.1 Scalar Field Definition}
Let $\phi(x)$ represent the coherent vibrational amplitude of filament bundles at point $x$:
\[
\phi(x) = \left\langle \text{Amplitude of coherent filament oscillations} \right\rangle_{F(x)}
\]

\subsection*{2.2 Covariant Equation}
The scalar field obeys the emergent Klein--Gordon equation:
\[
\left( \Box + m^2 \right) \phi(x) = 0
\]
where $\Box = g^{\mu\nu} \nabla_\mu \nabla_\nu$ is the d'Alembertian derived from the emergent metric, and the effective mass is:
\[
m^2 \sim \text{Topological strain energy density}.
\]

\section*{3. Path Integral over Filament Configurations}

\subsection*{3.1 Partition Function}
The quantum amplitude over the ensemble of filament bundles is:
\[
Z = \int \mathcal{D}F \, e^{i S[F]}
\]
where $F$ denotes filament configuration space and $S[F]$ is the emergent action.

\subsection*{3.2 Emergent Action}
\[
S[F] = \int d^4x \, \mathcal{L}_\text{eff}(F, \partial_\mu F)
\]
with $\mathcal{L}_\text{eff}$ capturing vibrational, strain, and phase coherence effects.

\section*{4. Emergent General Relativity}

\subsection*{4.1 Metric}
\[
g_{\mu\nu}(x) = \left\langle v_\mu v_\nu \right\rangle_{F(x)}
\]

\subsection*{4.2 Einstein Field Equations}
\[
\left\langle R_{\mu\nu} - \frac{1}{2} g_{\mu\nu} R \right\rangle = \kappa \left\langle E_\text{vib}(x), L_\text{link}(x) \right\rangle
\]
where $R_{\mu\nu}$ is the Ricci tensor, $R$ is the Ricci scalar, and the RHS represents emergent stress-energy from filament vibrational and topological energy.

\section*{5. Electromagnetism and Spinor Fields}

\subsection*{5.1 Gauge Potential}
\[
A_\mu(x) = \left\langle \text{Phase gradient structure} \right\rangle_{F(x)}
\]

\subsection*{5.2 Field Strength Tensor}
\[
F_{\mu\nu} = \partial_\mu A_\nu - \partial_\nu A_\mu
\]

\subsection*{5.3 Spinor Field Definition}
\[
\psi(x) = \left\langle \text{Local twist state of filament bundles} \right\rangle_{F(x)}
\]

\subsection*{5.4 Dirac Equation}
\[
\left( i\gamma^\mu (\partial_\mu + ie A_\mu) - m \right) \psi(x) = 0
\]

\subsection*{5.5 Current}
\[
J^\mu(x) = e \bar{\psi}(x) \gamma^\mu \psi(x)
\]

\section*{6. Summary Table}

\begin{center}
\begin{tabular}{|l|l|}
\hline
\textbf{Standard Field Quantity} & \textbf{SAT Interpretation} \\
\hline
$\phi(x)$ & Coherent filament vibration amplitude \\
$A_\mu(x)$ & Emergent from phase gradients \\
$\psi(x)$ & Local twist state of bundles (spin-$\frac{1}{2}$) \\
$F_{\mu\nu}$ & Shear/twist tensor of bundles \\
$T_{\mu\nu}$ & Vibrational + linking energy density \\
$g_{\mu\nu}$ & Averaged tangent product $\langle v_\mu v_\nu \rangle$ \\
\hline
\end{tabular}
\end{center}

\end{document}
